\section{Tidal Embassy: A Generation and Play Trace}
\label{app:tidal}


We now follow \emph{Tidal Embassy of Lost Languages} from premise to play. Its
one-sentence premise describes a
floating embassy where linguists, smugglers, ambassadors, and memory divers
trade endangered languages as political assets. The example preserves
intermediate images and refinement attempts, allowing each generative stage to
be connected to the final playable world.

\paragraph{Step 1: Premise analysis and planning.}
The planning stage identifies the nouns that become state---embassy,
languages, translation rights, archives, smugglers, diplomats, and tides---and
the verbs that imply gameplay---trade, translate, forge, authenticate, rescue,
and reveal. It then proposes an institutional loop: access to archives produces
linguistic objects; translators and forgers can
translate or forge them; diplomats exchange rights and debts; customs officers
can inspect transfers; flood events threaten stored knowledge. The resulting
world plan connects entities, conflicts, resource cycles, and open questions.

Planning establishes what must exist and why it matters. Later stages determine
spatial layout, visual appearance, and resource quantities. Missing activity
loops are refined before visual generation; for example, translation rights
must have a source, a use, and a means of transfer.

\paragraph{Step 2: Semantic structure.}
Semantic generation expands the plan into rooms, roles, items, institutions,
knowledge, and conflict hooks with stable identities. In this world,
the graph contains eight rooms, including the Phoneme Exchange, Mnemonic
Garden, Embassy Customs and Quarantine, and Ballast Hold; 25 named main
characters; and 20 distinct items. Explicit relations connect each character to
a home, each possession to an owner, and each action to relevant role,
knowledge, location, or resource conditions.

The generation process also preserves asymmetry. Shen Mo is a Phonetic
Smuggler and Lexical Restorationist based in the Cartographer's Nook. His
inventory includes a translation ear-cuff, archive key, forged exchange ledger,
dive helmet, ritual ink, and a fragment of the Unspoken Treaty. His private
knowledge includes phonetic decay patterns and ventilation shafts linking
Customs to the Mnemonic Garden. Another character receives different property,
authority, and information. This uneven allocation is what later makes
persuasion, trade, secrecy, and coalition formation possible.

\paragraph{Step 3: Visual direction and cross-modal identity.}
The visual direction converts semantic identity into a shared canon:
Art Nouveau diplomatic architecture, floating platforms, brass infrastructure,
bioluminescent language media, water intrusion, and a melancholic ritual tone.
It also emits per-entity appearance anchors. Room prompts inherit the canon but
add function-specific materials; character prompts inherit world palette and
wardrobe rules but retain role silhouettes; props receive icons whose shapes
communicate their use. Shared style and entity-specific appearance jointly keep
rooms, inhabitants, and objects recognizable as parts of the same world.

The visual direction guides each image independently while maintaining common
architecture, materials, palette, and role silhouettes. A visual inconsistency
can therefore be refined locally while the semantic identity of the world
remains stable.

\begin{figure*}[t]
\centering
\includegraphics[width=0.98\textwidth]{figures/tidal_embassy_visual_elements.pdf}
\caption{Generated visual elements from Tidal Embassy. The composed map contains
eight generated floors (A--D show the map and three constituent rooms); E--G
show the directional character atlases; H is a generated object. Shared
visual motifs connect spatial, character, and object identity.}
\label{fig:tidal-visual-elements}
\end{figure*}

\paragraph{Step 4: Evolving world state.}
This stage determines what can change after the images are produced.
It creates wallets, inventories, property libraries, knowledge assets, directed
relationships, task threads, room occupancy, and world-event status. It also
defines invariants. Item quantities remain nonnegative; physical transfers
conserve quantity; private knowledge has an owner and an acquisition event; a
character occupies one reachable coordinate; relationship changes are directed.
The generated repertoire and open-proposal forms are linked to the same world state.

The evolving state connects generated descriptions through observable causes
and consequences. The treaty fragment in Shen Mo's inventory can later be inspected,
given away, stolen through an allowed rule, or used as an input to another
action. When ownership changes, the next observation shows the new owner.

\paragraph{Step 5: Integrated world representation.}
The integrated representation connects semantic, visual, spatial, and evolving
state. It supports cross-domain questions:
Does every home room exist? Does every main character have a visual asset and a
walkable placement? Is an item used by an action defined and obtainable? Does a
private fact mention an existing institution? Does a room's purpose agree with
its objects and activities? These dependencies guide local refinement.

If a character image is inconsistent, roles, rooms, and items remain fixed. If
an item relation is missing, only the connected content is regenerated. Figure
\ref{fig:tidal-repair} illustrates this process: two room images are refined
while the rest of Tidal Embassy remains unchanged.

\paragraph{Step 6: Visual generation.}
FLUX generates the map source, individual room floors, characters, and local
objects from the shared visual direction. Character images are arranged into
directional movement atlases, and room images are placed within the planned
regions and connected by traversable corridors. Figure
\ref{fig:tidal-visual-elements} shows the composed map alongside three source floors,
three character atlases, and an object.

Functional visibility connects visual quality to use: room imagery must cover
its playable area, characters must remain recognizable while moving, and
furniture must preserve navigable paths.

\paragraph{Step 7: Visual and spatial evaluation.}
The evaluation examines semantic completeness, visual identity, traversable
geometry, foreground coverage, directional character consistency, readable
room composition, and visibility during play. The Diplomat's Lounge required
three real generation attempts. Attempt 1 was classified as a poster card with
an edge blank-background ratio of 0.56. Attempt 2 removed the card appearance
but retained a blank-margin failure with a maximum blank-row ratio of 0.73.
Only the third, edge-to-edge playable floor was accepted.

\begin{figure*}[t]
\centering
\includegraphics[width=0.98\textwidth]{figures/localized_visual_refinement.pdf}
\caption{Localized visual refinement for the Diplomat's Lounge. The first image
has a large poster-like margin; the second retains an extensive blank boundary;
the third fills the playable floor while preserving the room's identity. Other
rooms, characters, and world state remain unchanged.}
\label{fig:tidal-repair}
\end{figure*}

In \agora, visual quality combines coherent identity, readable composition, and
fitness for navigation and interaction. Preserved refinement attempts make the
progression from attractive illustration to usable world space observable.

\paragraph{Step 8: Persistent play.}
The generated world begins from a fixed initial state and evolves independently
for each session. Humans and AI agents encounter the same rooms, inhabitants,
objects, actions, and consequences. Separate world identities keep concurrent
multiworld experiments isolated.

Figure~\ref{fig:embodied-character} connects character generation to embodied play.
A crop of the interaction capture shows the selected agent. A complementary
world view shows neighboring agents among generated furniture and room surfaces,
while the directional atlas records appearance across locomotion. The interactive trial covers
movement, messaging, item use, trade, continued play, and a model-mediated reply.

The 25-character set uses one canonical visual identity per inhabitant.
Front locomotion and side views preserve the same hair, clothing, body type,
and palette, so an inhabitant remains recognizable while moving. Native
$64\times64$ frames retain facial, clothing, and limb detail in the stored atlases.

Character evaluation covers single-person composition, silhouette readability,
palette contrast, directional identity, absence of text and detached effects,
and visibility against the room surface. Contact-sheet review additionally
examines the complete cast together, revealing cross-view identity changes and
visual outliers that are difficult to judge one frame at a time.

Successive review passes identified characters with problematic props,
enclosing effects, orientation, text, detached sparkles, or insufficient
contrast. These counts reflect different review stages rather than disjoint
character sets.
Localized regeneration produced a final set of 25 coherent, readable character
atlases, all visible during the interactive trial.

\begin{figure*}[t]
\centering
\includegraphics[width=0.91\textwidth]{figures/embodied_character_interaction.pdf}
\caption{Character identity in Tidal Embassy. (A) A crop of the interaction capture shows the selected character,
marked by an amber ring. (B) A complementary
world view situates agents among generated room surfaces and furniture. (C) The
directional atlas preserves clothing, palette, and silhouette across
movement. The interactive trial includes movement, messaging, item use, and
trade.}
\label{fig:embodied-character}
\end{figure*}

\paragraph{Step 9: Interaction-informed refinement.}
During play, agents and humans observe bounded state and choose either a
world-specific action or an open proposal. The World Coordinator evaluates
effects and enacts successful transitions, as detailed in Section~\ref{sec:method}. Interaction traces
also provide feedback upstream. An action that is never selected may be poorly
named or irrelevant; an unreachable character suggests placement repair; an
important object that never enters an exchange suggests a missing action;
frequent open proposals around one intention suggest a useful addition to the
visible repertoire. Observed life inside the world can therefore guide the next
round of world refinement.

